\section{A Simple Model of Linguistic and Semantic Structure}
\label{sec:model}

Motivated by the empirical findings in Section~2, we now introduce a minimal
model that captures the essential geometric structure of language model
representations.
The goal of this model is not to faithfully describe all aspects of neural
representations, but to isolate the key properties needed to explain why
averaging provides an effective and safe approximation to spectral spike
suppression.

\subsection{Token Representation Model}

Let $x_i \in \mathbb{R}^H$ denote the representation of token $i$ in a given
layer, and let
\[
X = \begin{bmatrix} x_1^\top \\ \vdots \\ x_N^\top \end{bmatrix} \in \mathbb{R}^{N \times H}
\]
be the matrix of token representations in a batch.
We model each token representation as the sum of two components:
\begin{equation}
\label{eq:model_decomposition}
x_i = \ell + s_i,
\end{equation}
where:
\begin{itemize}
  \item $\ell \in \mathbb{R}^H$ is a \emph{shared linguistic component}, common to
  all tokens in the batch;
  \item $s_i \in \mathbb{R}^H$ is a \emph{semantic component} that varies across
  tokens.
\end{itemize}

This decomposition reflects the empirical observation that linguistic structure
(grammar, syntax, shared functional patterns) is coherent across tokens, whereas
semantic content is diverse and context-dependent.

\subsection{Structural Assumptions}

The empirical results in Sections~2.2--2.4 justify the following assumptions.

\paragraph{Assumption 1 (Low-rank linguistic structure).}
The shared component $\ell$ lies in a low-dimensional subspace
$\mathcal{L} \subset \mathbb{R}^H$ of dimension $d_\ell \ll H$.
Equivalently, $\ell$ is well-approximated by the span of a small number of
leading right singular vectors of $X$.

\paragraph{Assumption 2 (Semantic compressibility).}
Each semantic component $s_i$ lies predominantly in the orthogonal complement
$\mathcal{L}^\perp$, and is \emph{compressible} in a fixed semantic eigenbasis:
most of its energy is concentrated on a subset of coordinates whose size is
small relative to $\dim(\mathcal{L}^\perp)$.

\paragraph{Assumption 3 (Semantic incoherence across sequences).}
For tokens drawn from different sequences (or sufficiently distant positions in
long contexts), semantic residuals are weakly correlated in a \emph{normalized}
sense.
Let $s_i^{\perp} \triangleq P_{\mathcal{L}^\perp}s_i$ denote the projection of
$s_i$ onto the semantic (tail) subspace.
Then for $i\neq j$ from different sequences,
\begin{equation}
\label{eq:assump_cos_incoh}
\mathbb{E}\!\left[
\frac{\langle s_i^{\perp},\, s_j^{\perp}\rangle}
{\|s_i^{\perp}\|_2\,\|s_j^{\perp}\|_2}
\right] \approx 0,
\end{equation}
where the expectation is taken over batch sampling.
Empirically, this near-zero mean holds for cosine similarity of tail-basis
coefficients across sequences (Section~2.4).
In addition, semantic energy is bounded:
\begin{equation}
\label{eq:assump_sem_moment}
\mathbb{E}\|s_i\|_2^2 \le \sigma^2 < \infty.
\end{equation}
This assumption captures semantic diversity under batch sampling: different
sequences may activate overlapping semantic directions, but their coefficients
in the tail subspace are approximately unaligned, enabling cancellation under
averaging.

\subsection{Mean and Covariance Structure}

Let
\[
\mu = \frac{1}{N} \sum_{i=1}^N x_i
\]
denote the batch mean.
Substituting~\eqref{eq:model_decomposition}, we obtain
\[
\mu = \ell + \frac{1}{N} \sum_{i=1}^N s_i.
\]
By Assumption~3, semantic terms cancel in expectation under averaging, so $\mu$
is dominated by the shared linguistic component $\ell$.
This explains why the batch mean lies almost entirely in the spike subspace
(Section~2.3), and is quantified in Figure~\ref{fig:avg_mean_in_spike}.

Next, consider the centered representations
\[
\tilde{x}_i = x_i - \mu = s_i - \frac{1}{N} \sum_{j=1}^N s_j.
\]
The covariance of centered representations is therefore governed by the semantic
components:
\[
\frac{1}{N} \sum_{i=1}^N \tilde{x}_i \tilde{x}_i^\top
\;\approx\;
\frac{1}{N} \sum_{i=1}^N s_i s_i^\top,
\]
up to cross terms that vanish under Assumption~3.
This separation clarifies where information relevant to learning resides:
linguistic structure primarily manifests in shared, low-order statistics
(mean-like components), while semantic structure is encoded in second-order
variation (covariance and tail geometry).

\subsection{Averaging as Conservative Spike Suppression}

Define the mean-removed matrix
\[
X_{\mathrm{avg}} = X - \mathbf{1}\mu^\top,
\]
where $\mathbf{1} \in \mathbb{R}^N$ is the all-ones vector.
This operation removes a single rank-one component aligned with the batch mean.

Because $\mu$ lies in the low-rank linguistic subspace $\mathcal{L}$, mean removal
suppresses a direction that is contained in the spectral spike.
At the same time, because semantic components are incoherent across sequences,
mean removal has negligible effect on the semantic covariance.
Empirically, this behavior is reflected in two complementary ways:
(i) mean removal reduces spike energy (Figure~\ref{fig:avg_spike_reduction}),
and (ii) mean removal preserves tail energy almost entirely across a wide range
of spike-subspace definitions (Figure~\ref{fig:avg_tail_preservation}).

Importantly, averaging is \emph{conservative}:
it removes only one coherent direction rather than the entire spike subspace.
This is in contrast to SVD-based spike suppression, which removes the full
top-$k$ subspace.
Figure~\ref{fig:avg_residual_energy} illustrates this distinction: mean removal
removes a fixed fraction of total energy, while top-$k$ SVD removal continues to
remove energy as $k$ grows.

\begin{figure}[t]
  \centering
  \begin{subfigure}[t]{0.48\linewidth}
    \centering
    \includegraphics[width=\linewidth]{figures/avg_mean_in_spike_vs_k.png}
    \caption{Mean captured by spike subspace.}
    \label{fig:avg_mean_in_spike}
  \end{subfigure}
  \hfill
  \begin{subfigure}[t]{0.48\linewidth}
    \centering
    \includegraphics[width=\linewidth]{figures/avg_spike_reduction_vs_k.png}
    \caption{Spike energy reduction by mean removal.}
    \label{fig:avg_spike_reduction}
  \end{subfigure}

  \vspace{0.5em}

  \begin{subfigure}[t]{0.48\linewidth}
    \centering
    \includegraphics[width=\linewidth]{figures/avg_tail_preservation_vs_k.png}
    \caption{Tail energy preservation under mean removal.}
    \label{fig:avg_tail_preservation}
  \end{subfigure}
  \hfill
  \begin{subfigure}[t]{0.48\linewidth}
    \centering
    \includegraphics[width=\linewidth]{figures/avg_residual_energy_comparison.png}
    \caption{Residual energy: mean removal vs top-$k$ SVD removal.}
    \label{fig:avg_residual_energy}
  \end{subfigure}

  \caption{Quantitative validation of averaging as conservative spike suppression.
  (\textbf{Top-left}) The batch mean direction lies almost entirely in the spike
  subspace spanned by the leading right singular vectors.
  (\textbf{Top-right}) Mean removal reduces energy in the spike subspace, with
  strongest relative impact when the spike subspace is small.
  (\textbf{Bottom-left}) Tail energy is preserved to near unity under mean
  removal, indicating minimal disturbance of semantic variation.
  (\textbf{Bottom-right}) Mean removal is conservative compared to top-$k$ SVD
  removal: it removes a fixed shared component, whereas SVD removes increasing
  energy as $k$ grows.}
  \label{fig:avg_vs_svd_validation}
\end{figure}

\subsection{Empirical Effect of Averaging on Spike and Tail}
\label{subsec:avg_empirical}

We now empirically examine how mean removal (averaging) affects the spectral
structure identified in Section~\ref{subsec:spike_long_tail}.
Our goal in this subsection is purely observational: to quantify how much of the
spectral spike is suppressed by averaging, and how much of the long tail is
preserved, without appealing to any theoretical explanation.
A formal analysis of these phenomena is deferred to Section~4.

\paragraph{Experimental setup.}
Given the token matrix $X \in \mathbb{R}^{(BL)\times H}$, we consider two
operations:
(i) \emph{mean removal}, producing
$X_{\mathrm{avg}} = X - \mathbf{1}\mu^\top$, where $\mu$ is the batch mean; and
(ii) \emph{top-$k$ SVD removal}, producing
$X_k = X - X P_k$, where $P_k$ projects onto the top-$k$ right singular vectors
of $X$.
We evaluate these operations using Frobenius-norm energy ratios, which directly
measure how much variance remains after each transformation.

\paragraph{Residual energy under averaging and SVD removal.}
Figure~\ref{fig:avg_vs_svd_residual} compares the fraction of residual energy
$\|X_{\mathrm{res}}\|_F^2 / \|X\|_F^2$ after mean removal and after top-$k$ SVD
removal for varying $k$.
Mean removal suppresses a substantial portion of the spike energy, while leaving
the majority of total energy intact.
In contrast, increasing $k$ in SVD removal progressively removes more energy,
eventually eliminating both spike and tail components.
This comparison highlights that averaging acts as a conservative intervention,
removing less total energy than explicit low-rank projection.

\paragraph{Mean alignment with the spike subspace.}
To quantify how closely mean removal targets the spike, we measure the fraction
of mean energy captured by the spike subspace:
$\|P_k \mu\|_2^2 / \|\mu\|_2^2$.
As shown in Figure~\ref{fig:avg_mean_in_spike}, the batch mean lies almost entirely
within the spike subspace for modest values of $k$.
For example, projecting onto the top $1\%$ of singular directions captures more
than $99.9\%$ of the mean energy.
This observation indicates that mean removal primarily removes variance aligned
with the dominant spectral spike.

\paragraph{Preservation of tail energy.}
We next examine how averaging affects the long tail.
Figure~\ref{fig:avg_tail_preservation} reports the fraction of tail energy
preserved after mean removal as a function of the spike subspace size $k$.
Across a wide range of $k$, averaging preserves nearly all tail energy, even when
the spike subspace grows.
This behavior contrasts with top-$k$ SVD removal, which increasingly suppresses
tail energy as $k$ increases.

\paragraph{Summary of empirical findings.}
Taken together, these results show that mean removal:
(i) removes a nontrivial portion of the spectral spike;
(ii) targets directions closely aligned with the spike subspace; and
(iii) preserves the vast majority of long-tail energy.
These empirical regularities are robust across layers and batches in the models
we studied.

In the next section, we provide a formal explanation for why averaging exhibits
these properties, and show how they follow from the structural assumptions
validated in Sections~\ref{subsec:spike_coherent} and
\ref{subsec:tail_incoherent}.